\section{Discussion}
Reusing predictions changes the link between classification error and a
reported estimate. The relevant weights depend on the comparison, controls,
and regression role. Reuse can link errors across records and across
constructed variables even when distinct item errors are independent.
An outcome coefficient is linear in those errors, while a constructed
covariate changes the regression's cross-products and denominator.

This construction connects familiar inference problems. Generated-variable
methods address coefficient centering and estimated stages; shared-source
methods address dependence and influential loadings; validation and
survey methods address missing verified information. Their assumptions
must be checked for the same item-to-record map. A larger regression sample
does not itself supply more independently labeled or verified items.

For a probability audit, item and item-pair contributions give unbiased
regression moments and an unbiased joint covariance estimate. The normal
coefficient interval is a further approximation. The simulations show
substantial reductions in coefficient bias but also undercoverage in
concentrated and few-item designs. An exact covariance formula cannot
remove a weak denominator or create a central limit theorem. Reporting
these distinctions makes clear which correction has been justified and
which interval interpretation remains model dependent.

Confusion counts and ordinary score calibration preserve different
information. Counts generally omit the error locations needed by a weighted
estimate. Calibrated conditional means can support population plug-ins under
the relevant regression assumptions, but calibration averaged over items
does not establish those assumptions or the joint error covariance.
The validation design and retained item identities therefore determine
which downstream adjustments and uncertainty statements the data support.
